\documentclass[11pt,a4paper]{article}
\usepackage[T1]{fontenc}
\usepackage{lmodern,amsmath,amssymb,amsthm,mathtools,booktabs,microtype}
\usepackage[margin=26mm]{geometry}
\usepackage[colorlinks=true,linkcolor=blue,citecolor=blue,urlcolor=blue]{hyperref}
\setlength{\emergencystretch}{2em}
\newtheorem{theorem}{Theorem}[section]
\newtheorem{lemma}[theorem]{Lemma}
\newtheorem{proposition}[theorem]{Proposition}
\newcommand{\Q}{\mathbb Q}\newcommand{\Z}{\mathbb Z}
\DeclareMathOperator{\Tr}{Tr}\DeclareMathOperator{\Nm}{Nm}
\title{A quartic in 74 variables with small integer coefficients}
\author{}\date{16 November 2026}
\begin{document}\maketitle
\begin{abstract}
Under ordinary rational $abc$ and the stated elliptic, prime-pattern,
trace-height and simultaneous-twist premises, we construct a defining
quartic for $\Z$ in $\Q$ with 74 total variables. A complete quadratic
reference graph uses $R_0$ as its translating point and
exposes both reciprocal-abscissa squares. Its coefficients lie in
$\Z[\sqrt2]$. The full polynomial is a sum of 66 primitive integral
quadratic squares, has 1695 coordinate monomials and integer content one,
and its exact coefficient height is $\max\{30568,B^2+1\}$. The fixed
odd shift $B$ must be sufficiently large; its sufficient threshold is
ineffective. The source includes every torsion and tangent branch of
the retained paired chart and has a separate zero-input witness.
\end{abstract}

\section{Statement and arithmetic premises}
All variable counts include the input. For a reduced rational $u/v$,
$v>0$, put $h(u/v)=\log\max(|u|,v)$ and $h(0)=0$.
Set
\[
 K=\Q(\rho),\quad\rho^2=2,\quad D=11-6\rho,\quad
 E:y^2=x^3-Dx,\qquad R_0=(1+\rho,2\rho).
\]
Retain the arithmetic premises and their stated scope from
\cite{Parent}: $E(K)=\Z R_0+E[2](K)$ with $R_0$ non-torsion;
the fixed-data simultaneous-prime theorem with finite avoidance,
prescribed Gaussian congruences and local and real approximation;
the all-point trace-index bound; and the simultaneous congruent-twist
statement recalled in Section~\ref{sec:pole}. Ordinary rational $abc$
is the additional conjecture. The reference translator is $R_0$ itself. Write $m(Q)\in\Z$ for the
integer index in $Q=[m(Q)]R_0+T_0$, where $T_0\in E[2](K)$.
The complete reference and paired-point charts are algebraic results,
including their exceptional strata.
No stronger prime-pattern or simultaneous-twist assertion is used.

\begin{theorem}\label{thm:main}
Under these premises there is a computable family
\[
 Q_B\in\Z[X,Y_1,\ldots,Y_{73}],\qquad
 B\ge81\text{ a fixed odd integer},
\]
such that every member has degree exactly four and integer content one.
It is the sum of 66 primitive integral quadratic squares and has
1695 nonzero coordinate monomials. Its coefficient height is
\[
 H(Q_B)=\max\{30568,B^2+1\}.
\]
There is a fixed $B_*$ such that every odd $B\ge B_*$ gives rational
input projection exactly $\Z$. The shift $B$ is a coefficient, not a
quantified coordinate. No numerical sufficient value of $B_*$ is asserted.
\end{theorem}

The complete coefficient specification is the quadratic table in
\nolinkurl{computation/quartic/abc74_v48}. The field equations below,
the displayed affine substitutions and the stored invertible row
matrices uniquely determine that table. The value $B=81$ is an
algebraic test member, not a certified defining member. The exact
coefficient certificates do not constitute an independent audit of
every external arithmetic premise.

\section{The quadratic source}
Every field coordinate denotes two rational coordinates in the basis
$1,\rho$. A bar denotes the nontrivial automorphism of $K/\Q$.

Set
\[
 a_0=1+\rho,\qquad c_0=4-2\rho,\qquad
 \delta_0=-23+17\rho.
\]
Then $R_0=(a_0,a_0c_0)$ and $a_0\delta_0=D$.
Use five field coordinates $u_H,w_H,\ell,F_H,G_H$ and the affine
expressions
\[
 \begin{split}
 v_H&=a_0(u_H+w_H)+2c_0+\delta_0\ell,\\
 S&=a_0w_H+c_0+\delta_0\ell,\qquad
 T_H=-a_0u_H-c_0-\delta_0\ell.
 \end{split}
\]
Impose the five field quadrics
\begin{equation}\label{eq:reference}
\begin{aligned}
 a_0u_Hw_H&=1,\\
 \delta_0\ell^2+a_0\ell(u_H+w_H)+2c_0\ell&=1,\\
 a_0F_H(v_H+\delta_0w_H)+2c_0\ell u_H+\ell-u_H&=0,\\
 a_0G_H(v_H+\delta_0u_H)+2c_0\ell w_H+\ell-w_H&=0,\\
 a_0^2\{\delta_0^2F_HG_H-(1+4c_0w_H)F_H
                  -(1+4c_0u_H)G_H\}
       +1-2(2a_0+c_0^2)\ell^2&=0.
\end{aligned}
\end{equation}
These ten rational equations recover precisely
\[
 H\in E(K)\setminus\{O,(0,0),R_0,R_0+(0,0)\},
\]
with
\[
 \begin{gathered}
 f(H)=u_H\ell,\qquad f(H-R_0)=w_H\ell,\qquad f(Q)=1/x(Q),\\
 F_H=f(H)^2,\quad G_H=f(H-R_0)^2,\quad
 S=y(H)/x(H),\quad T_H=y(H-R_0)/x(H-R_0).
 \end{gathered}
\]
In particular $F_HG_H\ne0$ and $(S,T_H)\ne(0,0)$.
This is the complete inverse reference graph of~\cite[Theorem~2.1]{Reference}.
Its hypotheses hold because $a_0c_0\delta_0\ne0$ and
\[
 \Delta=(a_0+\delta_0)^2-4a_0c_0^2=1164-824\rho,
 \qquad \Nm(\Delta)=-3056\ne0.
\]
No square root of $a_0$ and no additional guard is required.
For clarity, its exceptional graph coefficients are also harmless.
The first two equations give $\ell v_H=1$ and $a_0u_Hw_H=1$.
The two graph coefficients $v_H+\delta_0w_H$ and
$v_H+\delta_0u_H$ cannot vanish together, since that would force
$\Delta=0$. If the first vanishes, the second equation for a square
already determines $G_H=1/\delta_0^2$, and the final equation has
nonzero $F_H$ coefficient $-4a_0^2c_0w_H$. The symmetric argument
handles the other case. Thus both squares are determined on every
reference fiber.

Use rational coordinates $X,\lambda,T,b_2,b_4,t_1$ and put
\[
                              \beta=X.
\]
The power and argument rows are
\begin{equation}\label{eq:powers}
 b_2=X^2,\qquad b_4=b_2^2,\qquad
 t_1-\lambda=B+b_4^2+T^2.
\end{equation}
Thus $b_2=\beta^2$, $b_4=\beta^4$, and the arguments are
$t_0=\lambda$, $t_1=\lambda+a$, where $a=B+\beta^8+T^2$.
The algebraic recovery arguments below assume $X\ne0$.
The zero input has explicit witnesses described in
Section~\ref{sec:completeness}.

For the point pair retain seven field coordinates
$f,h,g,k,s,F_U,F_V$ and two field ratios $L,L_0$, and impose
\begin{equation}\label{eq:pair}
 \begin{gathered}
 F_U=f^2,\qquad F_V=g^2,\qquad s(f-g)=h+k,\\
 s(h-k)=b_4-D(F_U+fg+F_V),
 \end{gathered}
\end{equation}
\begin{equation}\label{eq:ratios}
 \begin{aligned}
 hL+fS&=0,&(b_4-DF_U)L+hS&=0,\\
 kL_0+gT_H&=0,&(b_4-DF_V)L_0+kT_H&=0.
 \end{aligned}
\end{equation}
These rows are the complete paired chart after multiplying its
physical reciprocal abscissas by $\beta^2$, reciprocal ordinates by
$\beta^3$, chord parameter by $\beta$, and reciprocal-abscissa squares
by $\beta^4$, while dividing its two ratios by $\beta$.
For $X\ne0$, this substitution is invertible because $\beta\ne0$.
The paired-ratio theorem~\cite[Theorem~1.1]{Pair} recovers both
curve equations. The complete paired-chart theorem
\cite[Proposition~2.1]{Trace} then recovers
$U,V\in E(K)$, allowing identity numerators, with $U\ne V$ and
finite $A=U-V$. In particular
\[
 f=\beta^2/x(U),\quad g=\beta^2/x(V),\qquad
 x(A)=\beta^{-2}\bigl(s^2+D(f+g)\bigr).
\]
The physical ratios are
$L_{\rm phys}=\beta L=z(U)/z(H)$ and
$L_{0,\rm phys}=\beta L_0=z(V)/z(H-R_0)$ wherever these quotients are
defined, with $z=-x/y$. The complete equations also cover their
zero-ordinate strata.

Put $C=L-\overline{L_0}$, $I=L-1$. Impose
\begin{equation}\label{eq:common}
 C_2=C^2,\quad I_2=I^2,\qquad
                    b_2T=\Tr\bigl(s^2+D(f+g)\bigr).
\end{equation}
Thus $T=\Tr x(A)$ exactly, including finite torsion differences.

For each argument $t=t_j$ use rational coordinates
\[
        R,Z,w,\eta,b,e,u,p,v,q,p_2,q_2
\]
and field coordinates $A_1,A_2,\Omega=\omega_0+\rho\omega_1$.
Their contact rows are
\begin{equation}\label{eq:contactrows}
 R=t^2-2t,\quad RZ+Z-R^2+35R=0,\qquad
 (R+1)w=Z^2+24Z.
\end{equation}
Their detector and pole rows are
\begin{align}
 wA_1&=w^2+wC_2+F_HF_U,\label{eq:factor1}\\
 wA_2&=w^2+wI_2+\overline{G_H}\,\overline{F_V},\label{eq:factor2}\\
 w\Omega&=A_1A_2,\label{eq:product}\\
 \omega_0^2-2\omega_0\eta-2\omega_1^2&=0,
 \qquad wb=\eta^2,\label{eq:tracepole}\\
 p_2&=p^2,\qquad e^2+bp_2=1,\label{eq:pole1}\\
 uv&=e(1-b),\label{eq:pole2}\\
 q_2&=q^2,\qquad u^2-v^2=bq_2.\label{eq:pole3}
\end{align}
Every row is quadratic. The common part has 38 rational coordinates
and 34 rational rows. Each argument block has 18 coordinates and
16 rows. The totals are 74 coordinates and 66 rows.

\begin{lemma}\label{lem:nonzero}
Every rational solution with $X\ne0$ has $w\ne0$ in both blocks.
\end{lemma}
\begin{proof}
If $w=0$, the factor rows and $F_HG_H\ne0$ force
$F_U=F_V=0$, hence $f=g=0$. The complete paired chart then gives
$U=V=O$, contradicting $U\ne V$. Directly, the paired-ratio
identities first give the scaled reciprocal curve equations; these
force $h=k=0$, and the second row of~\eqref{eq:pair} contradicts
$b_4=\beta^4\ne0$. This includes $S=0$ or $T_H=0$, since they
cannot vanish simultaneously.
\end{proof}

\section{Local extraction with every torsion branch}
At good odd places, every finite valuation of a reciprocal abscissa
$f(Q)=1/x(Q)$ is even, with $f(O)=0$. Put $v(0)=+\infty$.

\begin{lemma}\label{lem:factor}
If $v(w)=\nu>0$ is odd and every finite valuation of $J$ is even,
then
\[
                  \mathcal D=1+C^2/w+J^2/w^2
\]
is nonzero and has valuation at most zero. If
$v(\mathcal D)>-\nu/2$, then $v(C)>\nu/4$ and
$v(J)>3\nu/4$.
\end{lemma}
\begin{proof}
The finite term valuations are $0$, $2v(C)-\nu$ and
$2v(J)-2\nu$. They lie respectively in the zero class modulo four,
an odd class, and the class two modulo four. Their minimum is unique.
The inequalities follow by comparing each nonconstant term with
$-\nu/2$. Zero summands cause no exception.
\end{proof}

\begin{lemma}\label{lem:branches}
Suppose a reference $H$ and numerator $U$ satisfy the complete
physical ratio equations, with $L_{\rm phys}=\beta L$ and $L$ a
local unit. If
\[
 v(\beta)\le0,\qquad v(\beta^2f(H)f(U))>0,
\]
then both points are formal at $O$.
\end{lemma}
\begin{proof}
The numerator cannot be $O$, because its regular ratio row would
force $L_{\rm phys}=0$. If $S=0$, the reference is a nonzero-abscissa
two-torsion point; the complete ratio rows with
$L_{\rm phys}\ne0$ force the numerator to be such a point too.
Both reciprocal abscissas are units at a good place, contradicting
the product inequality. Thus both ordinates are nonzero, and
$L_{\rm phys}=z(U)/z(H)$ is valid.

Put $a_0=v(f(H))$, $a_1=v(f(U))$,
$d_i=v(1-D f_i^2)$ and $t=v(\beta)$, where $f_0=f(H)$,
$f_1=f(U)$. The identity $z(Q)^2=f(Q)/(1-Df(Q)^2)$ gives
\[
 2t=a_1-d_1-a_0+d_0,\qquad
 v(\beta^2f(H)f(U))=2a_1+d_0-d_1.
\]
For $a_i>0$ one has $d_i=0$; for $a_i<0$, $d_i=2a_i$;
for $a_i=0$, $d_i\ge0$.
If $a_1<0$, positivity of the product gives $d_0>0$ and hence
$a_0=0$, so $2t=-a_1+d_0>0$, a contradiction.
If $a_1=0$, positivity gives $d_0>d_1\ge0$, so $a_0=0$ and
$2t=d_0-d_1>0$, again impossible.
Thus $a_1>0$ and $d_1=0$. If $a_0=0$, then
$2t=a_1+d_0>0$. If $a_0<0$, the positive product valuation is
$2(a_1+a_0)=4t$, also impossible. Therefore $a_0,a_1>0$,
which is exactly formal reduction at $O$.
\end{proof}

Let $n=m(A)\in\Z$ be the integer elliptic index of $A=U-V$
relative to $R_0$.
The contact rows give
\[
 w=\frac{\prod_{\theta\in\Theta}(t-\theta)}{(t-1)^6},
 \qquad\Theta=\{-5,-2,-1,0,2,3,4,7\}.
\]
Outside the fixed exceptional set, a consumed contact prime has
positive odd order $\nu=v_p(w)$.

\begin{proposition}\label{prop:local}
For $X\ne0$, at every consumed prime $p$, $\beta$ is $p$-integral and
\[
                         \beta(n-\beta)\equiv0\pmod p.
\]
\end{proposition}
\begin{proof}
The physical factors $\mathcal T_i=A_i/w$ are
\[
 \mathcal T_1=1+C^2/w+J_1^2/w^2,\quad
 \mathcal T_2=1+I^2/w+J_2^2/w^2,
\]
where $J_1=\beta^2f(H)f(U)$ and
$J_2=\beta^2\overline{f(H-R_0)}\,\overline{f(V)}$.
Each finite $J_i$ valuation is even. Lemma~\ref{lem:factor}
makes both factors nonzero of valuation at most zero at both places.
Write $W=\mathcal T_1\mathcal T_2$.
Equation~\eqref{eq:product} gives $\Omega=wW\ne0$.
If $\omega_0=0$, the trace row forces $\omega_1=0$, a contradiction.
Consequently
\begin{equation}\label{eq:physical}
 \tau=\frac{\Nm W}{\Tr W}\ne0,\qquad
 \eta=w\tau,\qquad b=w\tau^2\ne0.
\end{equation}
The native pole obstruction~\cite{Pole} excludes negative odd
valuation of $b$, so $v_p(\tau)>-\nu/2$.
The reciprocal-trace inequality
$v_p(\tau)\le\max_{\mathfrak p\mid p}v_{\mathfrak p}(W)$
selects one embedding where both factor valuations exceed $-\nu/2$.
Lemma~\ref{lem:factor} gives positive valuations of $C,I,J_1,J_2$.
Thus $L$ and $\overline{L_0}$ are units congruent to one, and both
physical ratios have valuation $t=v_p(\beta)$.

If $t>0$, the assertion follows. If $t\le0$,
Lemma~\ref{lem:branches} puts both reference/numerator pairs in their
formal groups at the indicated conjugate places. The reference labels
are $h_0,h_0-1\in\Z$, at least one a $p$-adic unit.
At that reference, the formal-logarithm ratio is integral, because
the numerator label is $p$-integral. Hence $t\ge0$, and therefore
$t=0$.

Both physical ratios now reduce to $\beta$. The formal logarithm
is an isometry on the formal group and agrees with its parameter
modulo positive order. It yields
\[
 m(U)\equiv h_0\beta,\qquad
 m(V)\equiv(h_0-1)\beta\pmod p.
\]
For a divisible reference label, integrality of the raw ratio forces
the numerator label to be divisible as well, giving the same congruence.
Subtracting proves $n\equiv\beta$. The argument includes all global
torsion labels and works at split and inert good primes.
\end{proof}

\section{Height soundness}
If $X=0$, the input is already an integer. Assume $X\ne0$.
By Lemma~\ref{lem:nonzero}, neither argument is a contact.
Ordinary rational $abc$ gives the eight-contact estimate
\cite[Section~2]{Height}
\[
                 h(t_j)\le\frac27\log M(t_j)+C_0.
\]
The mixed-power estimate at exponent $101/100$, applied to
$\beta^8+T^2$, gives
\begin{equation}\label{eq:height}
 h(\beta)\le\frac{301}{388}h(a)+C_1,\qquad
 h(T)\le\frac{301}{97}h(a)+C_2,\qquad h(a)\ge\log B.
\end{equation}
The marked-contact assertion uses ordinary $abc$; it is not asserted
from the single weak exponent alone. The all-point trace theorem
\cite[Corollary~1.3]{Trace}
supplies $m^2\le A_0h(T)+J_0$ for the integer index $m=n$
relative to $R_0$, after rescaling its fixed constants.

Suppose $n\ne\beta$. Write $\beta=c/d$ in lowest terms,
with $c\ne0$ and $d>0$, and put
\[
                         N=|c|\,|md-c|>0.
\]
Proposition~\ref{prop:local} gives $p\nmid d$ and $p\mid N$ at
every consumed prime. Hence both $M(t_j)$ divide $N$, and
\[
 h(a)\le\frac47\log N+C_3,\qquad
 N\le(1+|m|)e^{2h(\beta)}.
\]
Combining with~\eqref{eq:height} gives
\[
 h(\beta)\le\frac{602}{679}h(\beta)
       +\frac{301}{679}\log(1+|m|)+C_4.
\]
Strict absorption bounds $h(\beta)$, then $h(a)$ and $h(T)$,
by constants times $\log(1+|m|)+1$. The trace-index inequality
first bounds $|m|$, then all these heights, independently of $B$.
This contradicts $h(a)\ge\log B$ for all sufficiently large fixed
odd $B$.

Thus $n=\beta=X$, so $X=n=m\in\Z$. Together with the
zero-input case, this proves soundness on the entire rational zero set.

\section{Prepared pole inputs greater than one}\label{sec:pole}
The retained simultaneous-twist statement says: if $\beta_0>1$
and $\theta\ge1$ are coprime positive squarefree integers, both
$1\pmod8$, and each is a square modulo every prime of the other,
then some positive squarefree $r$ makes
$E_r:y^2=x^3-r^2x$ and $E_{r\theta}$ have infinite-order rational
points with abscissas in $\beta_0\Q^{\times2}$.
This is the retained application of Jaswal's
Theorem~1.2~\cite{Jaswal}, with the precise scope stated in~\cite{Pole}.

\begin{lemma}\label{lem:largepole}
Assume that statement. Suppose $b>1$ has square reduced denominator,
$b\in\Q_2^{\times2}$, $v_2(b)\le-4$, and every prime dividing
the squarefree part of its numerator is $1\pmod4$.
Then
\[
 e^2+bp^2=1,\qquad uv=e(1-b),\qquad u^2-v^2=bq^2
\]
has rational witnesses with $upq\ne0$.
\end{lemma}
\begin{proof}
Write $b=A/B_0^2$ reduced, with $A,B_0>0$, and write
$b=\beta_0 z^2$, $b-1=\theta c^2$, with $\beta_0,\theta$
positive squarefree and $z,c\ne0$. They are coprime, because
$\gcd(A,A-B_0^2)=1$.
The assumptions imply $4\mid B_0$, $A$ odd and $A\equiv1\pmod8$.
Thus both squarefree parts are $1\pmod8$.
At a prime dividing $\beta_0$, the residue of $b-1$ is $-1$;
it is a square because that prime is $1\pmod4$.
At a prime dividing $\theta$, the residue of $b$ is one.
Neither kind of prime divides $B_0$. These observations give the
two required mutual square-residue conditions.

If $\beta_0=1$, take $e=0$, $p=1/z$, $u=z$, $v=0$, $q=1$.
Otherwise apply the simultaneous-twist statement. Scaling its second
point by $(x,y)\mapsto(c^2x,c^3y)$ gives an infinite-order point
on $E_{r(b-1)}=E_{r(1-b)}$. Write the two points as
$(ba_1^2,y_1)$ on $E_r$ and $(ba_2^2,y_2)$ on $E_{r(1-b)}$,
with $a_1a_2y_1y_2\ne0$. Put
\[
 e=\frac{r}{ba_1^2},\quad p=\frac{y_1}{b^2a_1^3},\quad
 u=\frac{a_2}{a_1},\quad v=\frac{e(1-b)}u,\quad
 q=\frac{y_2}{b^2a_1^3u^2}.
\]
The first curve equation gives $bp^2=1-e^2$.
The second gives $bq^2=u^2-e^2(1-b)^2/u^2=u^2-v^2$.
The middle equation is the definition of $v$.
All displayed divisions are by nonzero quantities. The sign of
$1-b$ is retained in $v$; its square in the elliptic equation
justifies the indicated curve equality.
\end{proof}

\section{Completeness and preparation}\label{sec:completeness}
First fix $X\in\Z\setminus\{0\}$ and an odd $B\ge81$.
Before selecting any prime, set
\[
 \beta=X,\quad A=[X]R_0,\qquad
 T=\Tr x(A),\qquad a=B+\beta^8+T^2.
\]
The point $A$ is non-torsion for every nonzero integer $X$.
The contact lists $\Theta$ and $\Theta-a$ are disjoint.
Fix all denominator data, real neighborhoods and finite local
avoidance before applying the retained simultaneous-prime theorem.
In particular, exclude the numerator and denominator primes of the
nonzero abscissas and ordinates of $R_0$ and $A=[X]R_0$, at both embeddings,
as well as the primes dividing $2X$. Any further fixed nonzero
point coordinates needed for the ratio preparation may also be excluded.
All these data are fixed before prime selection.

The arbitrary-avoidance application in~\cite[Section~2]{Primes}
supplies $\lambda$, arbitrarily deep prescribed denominators, and
compatible order congruences for an integer $h_0$.
The selected Gaussian generators have real part $9$ modulo $16$
and imaginary part zero modulo $16$. The halving lemma
\cite[Section~5]{Half} puts $R_0$ and its conjugate in the doubled
subgroups of both reductions. Thus the order moduli
$\#E(\mathbf F_p)/4$ annihilate $R_0$.
The congruences $h_0\equiv0$ at the primary order and
$h_0\equiv1$ at the opposite order therefore suffice for
\[
 H=[h_0]R_0,\qquad U=[h_0X]R_0,\qquad
 V=[(h_0-1)X]R_0.
\]
This applies to odd and negative $X$ without an extra parity condition.

The same compact-group progression argument permits fixed real
neighborhoods avoiding the finitely many poles of these point and
ratio functions. Apply it to the finite family of multiples
$Q,XQ,X(Q-R_0)$ with $Q=[h_0]R_0$; different real components cause
no sign condition, since the detector terms are squares.
The finite local neighborhoods are also fixed before prime selection.
In particular $U-V=A$, so the previously fixed height data are unchanged.
The actual point functions define every weighted coordinate above.

At a selected simple contact, $\beta$ is a unit. At the primary
embedding put $Q=[h_0]R_0$, which is formal by the half-generator
argument. The formal multiplication expansion gives
\[
             \frac{z([X]Q)}{X z(Q)}
                         =1+O\bigl(z(Q)^4\bigr).
\]
The analogous formula holds for $Q'=[h_0-1]R_0$ at the opposite
reference embedding. Consequently the relative comparisons have
orders at least four and the relevant reciprocal abscissas have
orders at least two.
Hence the nonconstant terms of each factor have orders at least
seven and six. Each intended factor is a unit congruent to one.
At the opposite embedding the reference and numerator reduce to
$\pm R_0$ and $\pm A$. The prescribed abscissa avoidance makes
\emph{both} reciprocal abscissas units. The pair-product term has
order $-2$. The separation lemma makes each opposite factor have
negative order at most $-2$. Thus $W$ is a unit in one embedding
and has negative order in the other. Its trace cannot have
equal-order cancellation, and $v_p(\tau)=0$.
Therefore $v_p(b)=1$ at each selected prime.

The retained prime theorem, in the arbitrary-avoidance form
of~\cite{Primes}, selects every contact prime in the
class $1\pmod{8\cdot3\cdot17}$. More explicitly, with
$\lambda=u_0/L$, $a_0=0$, $a_1=a$ and all fixed denominators
absorbed in $L$, its positive prime values are
$p_{j,\theta}=u_0+(a_j-\theta)L$. The contact identity is
\begin{equation}\label{eq:squareclass}
 w_j=
 \frac{\prod_{\theta\in\Theta}p_{j,\theta}}
      {L^2\bigl(u_0+(a_j-1)L\bigr)^6}.
\end{equation}
The denominator is a rational square. Thus the squareclass of
$b_j=w_j\tau_j^2$ is supported only on the selected primes, all
$1\pmod4$. Their valuations are positive by the preceding local
calculation; every other valuation is even. The reduced denominator
of $b_j$ is therefore a square.

At the fixed preparation primes the reference, numerator and relative
ratio functions stay in fixed bounded local neighborhoods. Increasing
the prescribed denominator depth makes $1/w_j$ arbitrarily small.
The factors tend to one, $W_j$ tends to one and $\tau_j$ tends
to $1/2$. At two, the prescribed deep denominator gives
$v_2(w_j)$ negative even and, eventually, $v_2(\tau_j)=-1$.
Consequently $v_2(b_j)\le-4$. Formula~\eqref{eq:squareclass}
also gives $w_j\in\Q_2^{\times2}$, because all its numerator
primes are $1\pmod8$; hence $b_j$ is a two-adic square.

On the fixed real compact neighborhoods, as both arguments increase,
$w_j\sim t_j^2$, the factors tend uniformly to one and
\[
                          b_j=w_j\tau_j^2\sim t_j^2/4.
\]
The prescribed sufficiently large real argument therefore gives
$b_j>1$. Lemma~\ref{lem:largepole} applies to each actual physical
input, with independent pole witnesses. Equations
\eqref{eq:factor1}--\eqref{eq:tracepole} hold with
$A_i=w_j\mathcal T_i$, $\Omega=w_jW_j$ and $\eta=w_j\tau_j$.
The pole witnesses and their two squares satisfy the remaining rows.
This proves completeness for every nonzero integer.

For $X=0$, set $b_2=b_4=T=0$, set all seven scaled point
coordinates to zero, and put $L=L_0=1$, $C_2=I_2=0$.
Choose any allowed fixed reference point, for example $H=2R_0$.
Every point, ratio and trace row now holds literally.
Use the retained prime theorem for the two fixed contact lists with
$a=B$, taking sufficiently large real arguments and sufficiently deep
two-adic denominators. In each block set
\[
 A_1=A_2=w,\qquad \Omega=w,\qquad
 \eta=w/2,\qquad b=w/4.
\]
The detector rows hold, and formula~\eqref{eq:squareclass} gives
the same squareclass and square-denominator conclusions.
The prescribed two-adic depth and real lower bound give
$v_2(b)\le-4$ and $b>1$, while $b$ is a two-adic square.
Lemma~\ref{lem:largepole} supplies the remaining witnesses.
Thus the zero input also has a rational witness, without adding a
coordinate or a guard.

The physical input in this proof is $b=w\tau^2$ itself.
Its preparation is established directly by
\eqref{eq:squareclass}, the local factor orders and the two-adic and
real limits. No assertion about full-zero-set equivalence with
a differently normalized detector is needed.

\section{Exact coefficient specification and certificate}
Take the rational components of every displayed field equation,
using the basis $1,\rho$. Clear denominators, divide each row by its
integer content, and choose its sign so that the first nonzero
coefficient in the stored monomial order is positive. The coordinate
order is the 74-entry list in
\nolinkurl{computation/quartic/abc74_v48/quadratic-family.json}.

Apply the fixed $10\times10$ integer matrix in
\nolinkurl{basis-transformation.json} to the ten primitive reference
rows, and primitively normalize the resulting rows. Its determinant
is $\pm1$. This is a specified matrix operation; a lattice-reduction
algorithm is not part of the definition. Together with the other
56 primitive rows this gives rows $r_1,\ldots,r_{66}$. The first
row has constant coefficient one. For $i>1$, let $c_i$ be the
coefficient of the constant coordinate monomial and of $B^0$ in
$r_i$, and put
\[
 q_1=r_1,\qquad q_i=\operatorname{prim}(r_i-c_i r_1)
 \quad(2\le i\le66).
\]
Here $\operatorname{prim}$ is the same nonzero scalar normalization.
The parameter term $-B$ in the argument row is not included in $c_i$.
These row operations are invertible over $\Q$, so their common
rational zero set is exactly the full source just proved.

Define
\[
                         Q_B=\sum_{i=1}^{66}q_i^2.
\]
Over $\Q$, its zero set is exactly that common zero set. Exact
expansion gives 1695 nonzero coordinate monomials for every $B\ge81$
and 1697 parameter-separated terms. The input occurs only in
$X^2-b_2$, and no constant-elimination operation changes that row.
Hence the coefficient of $X^4$ is one, proving exact degree four
and integer content one for every integer $B$.

The constant coefficient is $B^2+1$. The largest absolute fixed
coefficient is 30568. The remaining parameter-dependent coefficients
have the form $c_0+c_1B$, with $|c_1|\le2$; the exact coefficient
certificate verifies their bounds by $B^2+1$ for $B\ge81$ and that
none vanishes in this range. This proves the asserted height formula
and support count. In particular, for odd $B\ge175$ the height is
exactly $B^2+1$; this coefficient crossover is not a sufficient
Diophantine defining threshold.

The producer reconstructs the literal field rows and the affine
reference outputs and verifies every stored matrix operation. A
separate checker uses rational-pair arithmetic, reconstructs all rows
without importing the producer, verifies the invertible basis changes,
and expands the squares independently. The retained complete paired
and pole rows are authenticated against their pinned exact source.
These finite checks certify the stated identities and coefficients;
they do not certify the external arithmetic premises.

\begin{thebibliography}{99}
\bibitem{Parent}\emph{A quartic in 74 variables from weighted point products},
\nolinkurl{papers/quartic_existence/quartic_abc74_20261116_v46.tex}.
\bibitem{Reference}\emph{Inverse reference functions and squared reciprocal abscissas},
\nolinkurl{papers/elliptic_charts/inverse_reference_squares_20261116_v48.tex}.
\bibitem{Pair}\emph{Recovery of elliptic curve equations from complete ratio equations},
\nolinkurl{papers/elliptic_charts/paired_ratio_redundancy_20261115_v40.tex}.
\bibitem{Trace}\emph{Trace heights on quadratic twists with $j=1728$},
\nolinkurl{papers/trace_heights/trace_on_all_points_20261113_v25.tex}.
\bibitem{Height}\emph{A quartic in 78 variables from a mixed-power height estimate},
\nolinkurl{papers/quartic_existence/quartic_abc78_20261115_v34.tex}.
\bibitem{Pole}\emph{Pole equations from simultaneous congruent-number twists},
\nolinkurl{papers/pole_parameters/congruent_twist_poles_20261114_v27.tex}.
\bibitem{Half}\emph{Heights of positive weighted elliptic traces},
\nolinkurl{papers/trace_heights/weighted_trace_20261114_v27.tex}.
\bibitem{Primes}\emph{Arithmetic dependencies of the construction with 29 variables},
\nolinkurl{papers/arithmetic_dependencies/arithmetic_dependencies_20261108_v8.tex},
Section~2.
\bibitem{Jaswal} M.~Jaswal, \emph{A proof that $\Z$ is Diophantine in $\Q$
through congruent-number twists}, Hexagon \texttt{2611.00331v2},
12 November 2026, Theorem~1.2, printed p.~4.
\url{https://hexagonmath.org/pdf/2611.00331v2}.
\end{thebibliography}
\end{document}
